\documentclass[12pt,a4paper]{article}
\usepackage[utf8]{inputenc}
\usepackage[T1]{fontenc}
\usepackage[polish]{babel}
\usepackage{amsmath}
\usepackage{amsfonts}
\usepackage{amssymb}
\usepackage{booktabs}
\usepackage{geometry}
\usepackage{enumitem}
\usepackage{parskip}
\geometry{margin=2cm}

\title{ZTBSK (2026L), 5. Transmisja dźwięku w sieciach komputerowych}
\author{Piotr Gugnowski, wydział EiTI, nr indeksu 320690}
\date{25 Maj 2026}

\begin{document}
\maketitle

\section*{Pytania i odpowiedzi}

% -----------------------------------------------------------------------
\subsection*{1. Jakie problemy związane są z transmisją dźwięku w sieciach
komputerowych?}

Transmisja dźwięku (VoIP) w sieciach opartych na protokole IP wiąże się
z szeregiem trudności:

\begin{itemize}
  \item \textbf{Brak gwarancji jakości w protokole IP.}
        IP jest protokołem datagramowym, niepewnym co do kolejności i dostarczenia
        pakietów. Protokół TCP zapewnia co prawda pewność dostarczenia, lecz
        kosztem opóźnień wynikających z potwierdzeń i ewentualnych retransmisji,
        które są niedopuszczalne przy transmisji głosu w czasie rzeczywistym.
  \item \textbf{Opóźnienia i ich kumulacja.}
        Głos wymaga dostarczenia pakietów w ściśle określonych granicach
        czasowych; opóźnienia wynikające z kodowania, kolejkowania i propagacji
        nie mogą się sumować powyżej dopuszczalnego progu (szczegóły w pytaniu 9).
  \item \textbf{Zmienność opóźnień (jitter).}
        Pakiety mogą docierać do odbiorcy w nierównych odstępach czasu, co
        zaburza odtwarzanie dźwięku. Konieczne jest stosowanie buforów dejitter.
  \item \textbf{Nieefektywne wykorzystanie pola ToS.}
        Pole \emph{Type of Service} nagłówka IPv4 definiuje klasy ruchu służące
        priorytetyzacji, lecz w praktyce większość urządzeń ich nie implementuje
        (szczegóły w pytaniu 7).
  \item \textbf{Niedobór adresów IP (IPv4).}
        Dołączenie dużej liczby telefonów IP do sieci może wymagać przenumerowania
        podsieci lub wymuszać rozbudowę infrastruktury routerów.
  \item \textbf{Współdzielenie łączy WAN z ruchem danych.}
        Łącza rozległe (szczególnie w sieciach Frame Relay i ATM) są
        współdzielone z ruchem danych, a ich przepustowość bywa ograniczona,
        co utrudnia zapewnienie wymaganej jakości transmisji głosu.
  \item \textbf{Blokowanie przez duże pakiety.}
        Duży pakiet danych wchodzący do kolejki może opóźnić pakiety głosowe
        na czas swojego wysyłania, co wymaga stosowania mechanizmów fragmentacji
        i kolejkowania.
\end{itemize}

% -----------------------------------------------------------------------
\subsection*{2. Jakie zagadnienia muszą zostać wzięte pod uwagę przy projektowaniu
multimedialnych sieci LAN, a jakie przy sieciach WAN?}

\textbf{Sieci LAN:}

W sieci lokalnej głównym zagadnieniem jest \textbf{priorytetyzacja ruchu
głosowego względem ruchu danych} przesyłanych tym samym łączem. Kwestie,
które należy wziąć pod uwagę:

\begin{itemize}
  \item \textbf{Kolejkowanie pakietów w przełącznikach.} Strumień głosowy
        (np.\ 80~kb/s per połączenie) współdzieli łącze z ruchem danych;
        przełącznik musi gromadzić pakiety danych tak, by nie opóźniać pakietów
        głosowych.
  \item \textbf{Klasyfikacja ruchu} (szczegóły w pytaniu 7). Standard
        IEEE~802.1p definiuje klasy ruchu w warstwie 2 (łącza danych);
        priorytety mogą być przenoszone do pola ToS nagłówka IP dla urządzeń
        warstwy 3 (sieciowej). Do celów telefonii wystarczają zazwyczaj dwie
        klasy: wysoka (głos) i niska (dane).
  \item \textbf{Adresowanie i segmentacja.} VLAN pozwala tworzyć logicznie
        odrębne sieci telefoniczne bez zmiany infrastruktury fizycznej.
        DHCP umożliwia automatyczną konfigurację telefonów IP, niezależnie
        od miejsca ich podłączenia.
\end{itemize}

\textbf{Sieci WAN:}

W sieciach rozległych dochodzą zagadnienia związane z \textbf{ograniczonym
i współdzielonym pasmem} łączy (Frame Relay, ATM) oraz większą podatnością
na opóźnienia. Kwestie do rozważenia:

\begin{itemize}
  \item \textbf{Algorytmy kolejkowania pakietów} (opisane szczegółowo
        w pytaniach 3 i 4): jawny podział pasma między klasy ruchu lub
        mechanizmy zapobiegania przeciążeniom.
  \item \textbf{Metody podniesienia wydajności łączy} (opisane w pytaniu 5):
        kompresja nagłówków, kompresja danych (kodeki), fragmentacja
        i przeplatanie pakietów (LFI).
  \item \textbf{Rezerwacja zasobów (RSVP)} (opisane w pytaniu 10): rezerwacja
        pasma na całej trasie pakietu gwarantująca end-to-end QoS.
  \item \textbf{Zarządzanie ruchem.} Listy routingu pozwalają kierować ruch
        głosowy dedykowanymi ścieżkami i wymuszać właściwe wartości pola ToS
        niezależnie od ustawień urządzeń końcowych.
\end{itemize}

% -----------------------------------------------------------------------
\subsection*{3. Do czego służy kolejkowanie pakietów?}

Kolejkowanie pakietów w urządzeniach sieciowych służy do \textbf{zarządzania
ruchem na wyjściu interfejsu} w sytuacji, gdy natężenie ruchu przychodzącego
przekracza chwilową przepustowość łącza. Podstawowe cele kolejkowania to:

\begin{itemize}
  \item Gromadzenie napływających pakietów danych tak, aby nie wprowadzać
        opóźnień w połączeniach głosowych (priorytetyzacja ruchu).
  \item Zapewnienie, że pakiety o wyższym priorytecie (np.\ głos) nie są
        wstrzymywane przez pakiety o niższym priorytecie (np.\ pobieranie plików).
  \item Zagwarantowanie każdej klasie ruchu określonej procentowo części
        dostępnego pasma łącza, niezależnie od aktualnego obciążenia.
  \item Zapobieganie przeciążeniom sieci przez wczesne wykrywanie i
        reagowanie na wzrost długości kolejek.
\end{itemize}

% -----------------------------------------------------------------------
\subsection*{4. Krótko opisz znane algorytmy działania kolejkowania pakietów
w urządzeniach sieciowych.}

\begin{description}
  \item[Priority Output Queuing (POQ)]
        Definiuje cztery kolejki per interfejs: z wysokim, normalnym, średnim
        i niskim priorytetem. Kolejki są opróżniane w tej kolejności; kolejka
        o niższym priorytecie obsługiwana jest tylko wtedy, gdy w kolejkach
        o wyższych priorytetach nie ma żadnych pakietów. Zaletą jest prostota
        i pewność, że ruch wysokopriorytetowy nie jest nigdy blokowany przez
        ruch niskoprzyorytetowy. Wadą jest brak gwarancji czasu dostępu do
        medium w każdej chwili (możliwe opóźnienie wynikające z pakietów
        aktualnie wysyłanych z niższej kolejki). Nadaje się do małych sieci.

  \item[Custom Queuing (CQ)]
        Umożliwia zdefiniowanie do 16 kolejek danych oraz jednej kolejki
        systemowej. Administratorzy przydzielają każdej kolejce procentową
        część dostępnego pasma interfejsu routera. Router cyklicznie sprawdza
        każdą kolejkę i wysyła z niej tyle bajtów, ile wynika z przydzielonego
        pasma. Gwarantuje to, że zarezerwowane pasmo jest zawsze dostępne dla
        danej klasy ruchu, nawet gdy łącze jest w pełni obciążone. Kolejki
        można definiować na podstawie protokołu, aplikacji lub numeru portu.
        Nadaje się dobrze do sieci LAN.

  \item[WRED (\emph{Weighted Random Early Detection})]
        Łączy mechanizm priorytetyzacji IP Precedence z algorytmem RED
        (\emph{Random Early Detection}). RED zapobiega przeciążeniom, wykrywając
        nadchodzące nasycenie kolejki i losowo odrzucając pakiety od nadawców
        generujących duży ruch, co wywołuje u nich zmniejszenie okna transmisji
        TCP. WRED dodaje do tego uwzględnienie priorytetów: pakiety z wyższym
        priorytetem są przepuszczane, a pakiety z niższym priorytetem są
        odrzucane wcześniej. Konfigurowane są progi i wagi dla pakietów
        o różnym priorytecie. Działa tylko dla rodziny TCP/IP; nadaje się
        do routerów wewnątrz sieci rozległych.
\end{description}

% -----------------------------------------------------------------------
\subsection*{5. Jakie metody mogą zostać zastosowane w celu podniesienia
wydajności łączy w transmisjach multimedialnych?}

\begin{description}
  \item[Fragmentacja i przeplatanie pakietów (LFI)]
        Duże pakiety danych są dzielone na mniejsze fragmenty i umieszczane
        w kolejce. Jeśli w tym czasie nadejdzie pakiet głosowy o wyższym
        priorytecie, zostaje przesunięty na początek kolejki i wysłany
        niezwłocznie. Samo kolejkowanie bez fragmentacji nie pomaga, bo duży
        pakiet wysyłany jest niepodzielnie; dopiero połączenie obu technik
        zapewnia niskie opóźnienia dla głosu.

  \item[Kompresja nagłówków]
        Sumaryczna wielkość nagłówków IP + UDP + RTP wynosi 40~bajtów, podczas
        gdy same dane głosowe to zwykle 150~bajtów. Kompresja nagłówków
        zmniejsza ich łączny rozmiar do 2--5~bajtów. Ma to istotne znaczenie
        przy transmisji dużej liczby krótkich pakietów głosowych, szczególnie
        na łączach o małej szybkości.

  \item[Kompresja danych (kodeki głosowe)]
        Kodeki zmniejszają wymagane pasmo wielokrotnie względem nieskompresowanego
        PCM (64~kb/s), kosztem pewnego opóźnienia kodowania i nieznacznego
        obniżenia jakości. Przegląd kodeków i ich parametrów zawiera pytanie 8.

  \item[Listy routingu]
        Pozwalają przydzielić ruchowi głosowemu dedykowane interfejsy wyjściowe,
        adres następnego routera lub wymusić konkretną wartość pola ToS,
        zapewniając specjalne traktowanie połączeń głosowych przez całą sieć.
\end{description}

% -----------------------------------------------------------------------
\subsection*{6. Do czego służy protokół RTP?}

RTP (\emph{Real Time Protocol}, RFC~3550) jest protokołem warstwy aplikacji
przeznaczonym do \textbf{przesyłania danych multimedialnych w czasie rzeczywistym},
takich jak transmisja głosu i wideo. Działa nad protokołem UDP (a nie TCP),
by uniknąć opóźnień wynikających z mechanizmów niezawodności.

Protokół RTP zapewnia:
\begin{itemize}
  \item \textbf{Numerowanie sekwencyjne pakietów} (ang.\ \emph{sequence number}),
        umożliwiające odbiorniku wykrycie utraty lub zmiany kolejności pakietów.
  \item \textbf{Znaczniki czasu} (ang.\ \emph{timestamp}), pozwalające na
        prawidłową synchronizację odtwarzania strumienia i obliczanie jittera.
  \item \textbf{Identyfikację typu ładunku} (ang.\ \emph{payload type}),
        informującą o zastosowanym kodeku (np.\ G.711, G.729, H.264).
  \item \textbf{Identyfikator źródła} (SSRC), umożliwiający rozróżnienie wielu
        strumieni w ramach jednej sesji.
\end{itemize}

RTP współpracuje z protokołem RTCP (\emph{RTP Control Protocol}), który
dostarcza statystyki sesji: utratę pakietów, jitter i opóźnienie.

% -----------------------------------------------------------------------
\subsection*{7. Wyjaśnij zastosowanie pola ToS nagłówka IP w sieciach
wielousługowych.}

Pole \textbf{ToS} (\emph{Type of Service}) w nagłówku IPv4 służy do
\textbf{oznaczania pakietów priorytetem}, dzięki czemu urządzenia sieciowe
mogą różnicować sposób ich traktowania. W sieciach wielousługowych (głos,
wideo, dane) pozwala to zagwarantować niskie opóźnienia dla ruchu
wrażliwego czasowo przy jednoczesnym przepuszczaniu ruchu o niższym priorytecie
w miarę dostępności pasma.

W praktyce stosuje się mechanizm \textbf{IP Precedence}, który wykorzystuje
trzy pierwsze bity pola ToS do zdefiniowania czterech klas priorytetu
(dwie są zarezerwowane). Priorytety mogą być nadawane przez urządzenia końcowe
(telefony IP, oprogramowanie) lub routery brzegowe, a następnie respektowane
przez wszystkie urządzenia na trasie pakietu.

Kluczowym zagadnieniem przy wdrożeniu jest wiarygodność oznaczeń. Stosuje
się dwa modele:
\begin{itemize}
  \item \textbf{Zaufane połączenie graniczne} (\emph{trusted boundary}):
        administrator kontroluje urządzenia po obu stronach połączenia
        i ma pewność, że wartości ToS są ustawiane prawidłowo.
  \item \textbf{Wymuszona klasyfikacja na portach przełącznika}: dla połączeń
        z urządzeniami zewnętrznymi porty przełącznika są konfigurowane tak,
        by nadpisywać oznaczenia ToS zgodnie z lokalną polityką, niezależnie
        od ustawień urządzeń końcowych.
\end{itemize}

Pole ToS, wraz z protokołem RSVP (rezerwacja pasma end-to-end, zob.\ pytanie 10),
tworzy podstawę mechanizmów QoS na całej trasie pakietu.

% -----------------------------------------------------------------------
\subsection*{8. Czym różnią się między sobą metody kodowania głosu?}

Metody kodowania głosu różnią się przede wszystkim \textbf{wymaganym pasmem
transmisji}, \textbf{opóźnieniem wprowadzanym przez koder} oraz
\textbf{jakością odtwarzanego dźwięku} (mierzoną oceną MOS,
\emph{Mean Opinion Score}, w skali 0--5):

\begin{description}
  \item[PCM (G.711, \emph{Pulse Code Modulation})]
        Standardowa metoda stosowana w tradycyjnej telefonii. Próbkuje sygnał
        analogowy z częstotliwością 8~kHz (wynikającą z twierdzenia Nyquista
        dla pasma 200--3400~Hz), każda próbka kodowana jest 8 bitami, co daje
        przepływność 64~kb/s. Opóźnienie: 0,125~ms. Ocena MOS: 4,1.
  \item[AD-PCM (G.726, \emph{Adaptive Differential PCM})]
        Koduje różnicę między kolejnymi próbkami zamiast wartości bezwzględnej,
        uzyskując 32~kb/s przy zachowaniu jakości porównywalnej z PCM.
        Opóźnienie: 1~ms. Ocena MOS: 4,1.
  \item[CS-ACELP (G.729, \emph{Code-Excited Linear Prediction})]
        Metoda parametryczna: nie przesyła spróbkowanego sygnału, lecz jedynie
        parametry modelu głosu. Wymagane pasmo: 8~kb/s. Opóźnienie: 10~ms.
        Ocena MOS: ok.\ 3,92.
\end{description}

Ogólna zasada jest następująca: im bardziej zaawansowana metoda kompresji
(niższe pasmo), tym większe opóźnienie kodowania i nieco niższa jakość. Ocena
MOS poniżej 3,0 wyklucza praktyczne zastosowanie do rozmów telefonicznych.
Wpływ opóźnienia kodowania na łączne opóźnienie transmisji omawia pytanie 9.

% -----------------------------------------------------------------------
\subsection*{9. Jakie czynniki wpływają na powstawanie opóźnień w transmisji
pakietów? Dlaczego te opóźnienia mają znaczenie przy transmisji głosu?}

\textbf{Źródła opóźnień:}

\begin{itemize}
  \item \textbf{Opóźnienie kodowania i kompresji.}
        Koder musi zebrać próbki dźwięku zanim wyśle pakiet. Dla G.729:
        co 10~ms generowany jest nowy pakiet, a jego kompresja zajmuje kolejne
        15~ms, więc pierwszy pakiet jest wysyłany dopiero po 25~ms.
  \item \textbf{Opóźnienie kolejkowania w urządzeniach sieciowych.}
        Pakiety czekają w kolejkach routerów i przełączników. Na łączu 64~kb/s
        pakiet 1500~bajtów może być transmitowany przez ok.\ 188~ms.
  \item \textbf{Opóźnienie propagacji.}
        Czas przejścia sygnału przez medium fizyczne (ok.\ 5~ms/1000~km dla
        kabla; znacznie więcej dla satelity).
  \item \textbf{Opóźnienie przetwarzania w routerach.}
        Każdy router dodaje kilka do kilkunastu milisekund.
  \item \textbf{Duże pakiety blokujące dostęp do łącza.}
        Router wysyłający duży pakiet danych nie może jednocześnie wysłać
        pakietu głosowego; bez fragmentacji tworzy się opóźnienie równe
        czasowi transmisji całego dużego pakietu.
\end{itemize}

\textbf{Znaczenie opóźnień dla transmisji głosu:}

Transmisja głosu ma charakter interaktywny i odbywa się w czasie rzeczywistym.
Opóźnienie powyżej 150--200~ms (ITU-T G.114) jest wyraźnie odczuwalne przez
rozmówców jako irytujące echo lub nakładanie się wypowiedzi. W odróżnieniu od
transferu plików, przy transmisji głosu nie można po prostu buforować i
retransmitować opóźnionych pakietów: spóźniony pakiet jest bezużyteczny
i należy go odrzucić. Dlatego wszystkie opóźnienia muszą mieścić się w ścisłych
granicach czasowych, a ich minimalizacja jest kluczowym celem projektowania
sieci VoIP.

% -----------------------------------------------------------------------
\subsection*{10. Do czego służy protokół RSVP?}

RSVP (\emph{Resource Reservation Protocol}, RFC~2205) jest protokołem
sygnalizacyjnym służącym do \textbf{rezerwacji zasobów sieciowych} (przede
wszystkim pasma i buforów) na całej trasie pakietu między nadawcą a odbiorcą.

Działanie: aplikacja żądająca transmisji w czasie rzeczywistym (np.\ rozmowa
VoIP) wysyła żądanie rezerwacji zasobów przez wszystkie routery na ścieżce do
odbiorcy. Każdy router, który akceptuje żądanie, rezerwuje określone pasmo
dla danego strumienia danych i odrzuca żądanie, jeśli zasoby są niedostępne.

Zastosowanie RSVP razem z polem IP Precedence (lub polem DSCP, \emph{Differentiated
Services Code Point}) pozwala zapewnić \textbf{end-to-end QoS} na całej
trasie pakietu, również przez sieci rozległe oparte na różnych protokołach
warstwy pierwszej i drugiej (Frame Relay, ATM, SONET, PPP). Dzięki temu
możliwe jest zagwarantowanie parametrów transmisji (pasmo, opóźnienie, jitter)
wymaganych przez aplikacje multimedialne działające w czasie rzeczywistym.

\end{document}
